\documentclass{report}
\usepackage{amsmath,amssymb}
\setlength{\parindent}{0mm}
\setlength{\parskip}{1em}
\begin{document}
\begin{center}
\rule{6in}{1pt} \
{\large Jennifer Pestana \\
{\bf Combination preconditioning of saddle point systems for positive definiteness }}

Mathematical Institute \\ University of Oxford \\ 24-29 St Giles\’ \\ Oxford OX1 3LB \\ United Kingdom
\\
{\tt pestana@maths.ox.ac.uk}\end{center}

Preconditioned iterative methods in nonstandard inner products for saddle
point systems have recently received attention. Krzy\.zanowski
(\emph{Numer. Linear Algebra Appl.} 2011; {\bf18}:123--140) identified a
two-parameter family of preconditioners in this context and Stoll and
Wathen (\emph{SIAM J. Matrix Anal. Appl.} 2008; {\bf 30}:582--608)
proposed combination preconditioning, where two different
preconditioners---for each of which the preconditioned saddle point
matrix is self-adjoint with respect to an inner product---can be blended
to create additional preconditioners and associated bilinear forms or
inner products. If a preconditioned saddle point matrix is nonsymmetric
but self-adjoint with respect to a nonstandard inner product a
MINRES-type method ($\cal W$-PMINRES) can be applied in the relevant
inner product. If the preconditioned matrix is also positive definite
with respect to this inner product a more efficient CG-like method ($\cal
W$-PCG) can be used reliably.

In this talk we provide explicit expressions for the combination of
certain Krzy\.{z}anowski preconditioners. We additionally give an example
of the rather counterintuitive result that the combination of two
preconditioners for which only $\cal W$-PMINRES can be reliably used can
lead to a preconditioner for which, for certain parameter choices, $\cal
W$-PCG is applicable. That is, the resulting preconditioned saddle point
matrix is positive definite with respect to an inner product. We show
that this combination preconditioner outperforms either of the two
preconditioners from which it is formed for a number of test problems.


\end{document}
